% !PRINTED 43
\weberpage{43}

Among the various properties of the binomial coefficients we shall derive two, of which we shall afterwards make an interesting application.

From (1), replacing $x,y$ by $1,x$ and taking $n=0,1,2,3\ldots$, the formulas
\begin{equation}
\begin{aligned}
1&=B_0^{(0)},\\
1+x&=B_0^{(1)}+B_1^{(1)}x,\\
(1+x)^2&=B_0^{(2)}+B_1^{(2)}x+B_2^{(2)}x^2,\\
&\cdots\\
(1+x)^n&=B_0^{(n)}+B_1^{(n)}x+B_2^{(n)}x^2+\cdots+B_n^{(n)}x^n.
\end{aligned}
\tag{2}
\end{equation}
are obtained.

We now use the well-known summation formula for the geometric series:
\[
1+(1+x)+(1+x)^2+\cdots+(1+x)^n={(1+x)^{n+1}-1\over x},
\]
where, if $(1+x)^{n+1}$ is again expanded by the binomial formula, replacing $n$ by $n+1$ in the last formula (2) and setting $B_0^{(n+1)}=1$,
\begin{equation}
{(1+x)^{n+1}-1\over x}=B_1^{(n+1)}+B_2^{(n+1)}x+\cdots+B_{n+1}^{(n+1)}x^n.
\tag{3}
\end{equation}
Comparing this with the sum of the right hand sides of (2), and, according to the rules of \S~1, setting equal the coefficients of corresponding powers of $x$, one obtains
\begin{equation}
\begin{aligned}
 B_0^{(0)}+B_0^{(1)}+B_0^{(2)}+\cdots+B_0^{(n)}&=B_1^{(n+1)},\\
 B_1^{(1)}+B_1^{(2)}+\cdots+B_1^{(n)}&=B_2^{(n+1)},\\
&\cdots\\
 B_n^{(n)}&=B_{n+1}^{(n+1)},
\end{aligned}
\tag{4}
\end{equation}
or in general
\begin{equation}
 B_\nu^{(\nu)}+B_\nu^{(\nu+1)}+\cdots+B_\nu^{(n)}=B_{\nu+1}^{(n+1)}.
\tag{5}
\end{equation}
If however the equations (2) are multiplied in order by
\[
 B_0^{(n)},
 \quad -B_1^{(n)},
 \quad +B_2^{(n)},
 \ldots,
 \quad \pm B_n^{(n)}
\]
where the upper sign is to be taken for even and the lower for odd $n$, the sum of the left hand sides gives
\[
B_0^{(n)}-B_1^{(n)}(1+x)+B_2^{(n)}(1+x)^2-\cdots\pm B_n^{(n)}(1+x)^n=[1-(1+x)]^n=(-x)^n
\]
and equating the coefficients of equal powers on the right and left hand sides yields the system of formulas:
%CONT>>
